\section{Software: TRUST Platform}
\label{sec:WP7:TRUST Platform:software}

\begin{itemize}
    \item \textbf{Contact Email(s):} pierre.ledac@cea.fr
    \item \textbf{Supported Architecture(s):} HYBRID
    \item \textbf{Repository Link:} \href{https://github.com/cea-trust-platform}{https://github.com/cea-trust-platform}
\end{itemize}

\subsection{Software Overview}
\label{sec:WP7:TRUST Platform:summary}

TRUST is CEA's modular simulation platform for multi-physics thermal-hydraulics and neutronics applications. Under WP7 the team
is aligning TRUST with the shared CI/CD, packaging, and benchmarking practices: container images based on the WP7 Spack stack
are now produced automatically, and key demonstrators (two-phase flow, porous media) are being wrapped with ReFrame test suites.
The platform also acts as an integration point for WP2 surrogate models and WP6 uncertainty quantification pipelines.

\subsection{Parallel Capabilities}
\label{sec:WP7:TRUST Platform:performances}


TRUST combines MPI for domain decomposition with on-node OpenMP parallelism and optional Kokkos-based kernels for accelerator
experiments. Production deployments run on TGCC IRENE and CEA's internal hybrid clusters, scaling to more than 6000 CPU cores for
two-phase flow benchmarks. The platform interfaces with PETSc, Trilinos, and ParaView Catalyst; these dependencies are
harmonised with WP7's shared Spack recipes to guarantee interoperability with the rest of the Exa-MA software stack.

\subsection{Initial Performance Metrics}
\label{sec:WP7:TRUST Platform:metrics}

The initial WP7 benchmarking wave covers the single-phase thermal-hydraulics demonstrator ``app-trust-th'' with datasets shared
through the WP7 TRUST Zenodo community (publication pending). Automated ReFrame runs on IRENE's Rome nodes report 0.72~s/step at 1024 cores (89\%
weak-scaling efficiency relative to 128 cores). Energy traces confirm a 13\% reduction compared to legacy scripts once WP7's
NUMA pinning and I/O staging are applied. Remaining bottlenecks include long preprocessing phases in Python steering scripts and
limited GPU coverage; both are being profiled ahead of the next release cycle.



\subsubsection{Benchmark \#1}
\begin{itemize}
    \item \textbf{Description:} ``app-trust-th'' weak-scaling benchmark on IRENE (128--1024 cores) modelling turbulent flow in a heated channel.
    \item \textbf{Benchmarking Tools Used:} ReFrame orchestrator, Extrae MPI traces, WP7 energy plug-in for RAPL measurements, and Perf/TMA counters for hotspot identification.
    \item \textbf{Results Summary:} 89\% weak-scaling efficiency; 13\% energy reduction after WP7 NUMA pinning; memory footprint reduced by 9\% using compressed checkpoints.
    \item \textbf{Challenges Identified:} Steering scripts create serial start-up overhead; GPU porting limited to pilot kernels---work continues with WP3 optimisation experts.
\end{itemize}

\subsection{12 months roadmap}
\label{sec:WP7:TRUST Platform:roadmap}

\begin{itemize}
    \item \textbf{Data improvements:} Release a public subset of TRUST benchmark datasets (geometry + BCs) with reproducibility documentation aligned with WP7 FAIR guidelines (Q1~2026).
    \item \textbf{Methodology application:} Extend CI coverage to GPU-enabled builds (Kokkos back-end) and integrate nightly regression tests on hybrid partitions (Q2~2026).
    \item \textbf{Results retention:} Publish quarterly performance reports to the Exa-MA knowledge base with trend analysis for runtime, energy, and solver robustness metrics (starting Q3~2026).
\end{itemize}
